\section{Scoring contract}
\label{app:scoring}

This appendix gives the formal definitions of the metrics introduced in Section~\ref{sec:scoring}. The scoring code lives in \texttt{src/causal\_discovery/scoring/scores.py}; the submission object is defined in \texttt{src/causal\_discovery/scoring/submission.py}.

A submission is a partially directed graph $\mathcal{S} = (V_d, E_d, E_u)$ over the public node set, where $E_d$ is a directed edge set whose directed subgraph is itself a DAG and $E_u$ is a set of canonical undirected edges $(a, b)$ with $a < b$. The two edge sets are disjoint: a pair cannot be both directed and unresolved. Self-loops and 2-cycles are rejected at submission time. The skeleton of $\mathcal{S}$ is
\begin{equation}
\mathrm{skel}(\mathcal{S}) = \pi(E_d) \cup E_u,
\end{equation}
where $\pi$ projects each directed edge to its canonical unordered pair. All scores below operate on the public-labelled objects; the anonymisation permutation of Appendix~\ref{app:instance:anon} is applied symmetrically to truth and submission, so it cancels.

For a TP/FP/FN triple, precision, recall, and F1 follow the standard definitions
\begin{equation}
\mathrm{P} = \frac{\mathrm{TP}}{\mathrm{TP} + \mathrm{FP}}, \qquad
\mathrm{R} = \frac{\mathrm{TP}}{\mathrm{TP} + \mathrm{FN}}, \qquad
\mathrm{F1} = \frac{2 \, \mathrm{P} \cdot \mathrm{R}}{\mathrm{P} + \mathrm{R}},
\end{equation}
with the convention that $\mathrm{P} = \mathrm{R} = 1$ whenever the relevant denominator is zero, and $\mathrm{F1} = 0$ if $\mathrm{P} + \mathrm{R} = 0$.

\subsection{Skeleton F1}
\label{app:scoring:skeleton}

Skeleton F1 measures adjacency recovery against the observational ceiling. Let $\mathrm{skel}(\mathcal{C}_\pi)$ denote the skeleton of the true CPDAG (equivalently, the skeleton of the hidden DAG). The submitted skeleton is $\mathrm{skel}(\mathcal{S})$ defined above. Counts are over unordered pairs:
\begin{equation}
\mathrm{TP} = |\mathrm{skel}(\mathcal{S}) \cap \mathrm{skel}(\mathcal{C}_\pi)|, \quad
\mathrm{FP} = |\mathrm{skel}(\mathcal{S}) \setminus \mathrm{skel}(\mathcal{C}_\pi)|, \quad
\mathrm{FN} = |\mathrm{skel}(\mathcal{C}_\pi) \setminus \mathrm{skel}(\mathcal{S})|.
\end{equation}
Skeleton F1 is symmetric in directed and unresolved submitted edges: both project to the same skeleton.

\subsection{Compelled-edge F1}
\label{app:scoring:compelled}

Let $E^c = E_d(\mathcal{C}_\pi)$ be the compelled edges of the true CPDAG (its directed portion). Compelled-edge scoring iterates only over $E^c$. For each $(u, v) \in E^c$:
\begin{equation}
\begin{aligned}
\text{if } (u, v) \in E_d(\mathcal{S}) &\quad \Rightarrow \quad \mathrm{TP} \mathrel{+}= 1, \\
\text{else if } (v, u) \in E_d(\mathcal{S}) &\quad \Rightarrow \quad \mathrm{FP} \mathrel{+}= 1 \text{ and } \mathrm{FN} \mathrel{+}= 1, \\
\text{else} &\quad \Rightarrow \quad \mathrm{FN} \mathrel{+}= 1.
\end{aligned}
\end{equation}
Two consequences are worth flagging. First, a reversed compelled edge is penalised twice, once as a false positive and once as a false negative, because it both asserts an incorrect orientation and fails to assert the correct one. Second, submitted directed edges that do not project to any pair in $\pi(E^c)$ contribute nothing to the compelled-edge counts; they are scored elsewhere. This makes compelled-edge F1 a directed measure restricted to the observationally identifiable orientations, not a directed F1 evaluated on the CPDAG.\footnote{When $E^c = \emptyset$ (the CPDAG has no compelled edges), the convention $\mathrm{P} = \mathrm{R} = 1$ applies and compelled-edge F1 equals $1$ trivially. Such instances are excluded from the in-paper levels by construction (Appendix~\ref{app:instance:filters}: the rejection filters require at least two undirected edges and a non-empty minimum sufficient intervention set).}

\subsection{Directed F1}
\label{app:scoring:directed}

Directed F1 is the symmetric set-comparison of submitted directed edges against the hidden DAG:
\begin{equation}
\mathrm{TP} = |E_d(\mathcal{S}) \cap E(G_\pi)|, \quad
\mathrm{FP} = |E_d(\mathcal{S}) \setminus E(G_\pi)|, \quad
\mathrm{FN} = |E(G_\pi) \setminus E_d(\mathcal{S})|.
\end{equation}
A submitted undirected edge $\{u, v\} \in E_u$ contributes to none of these counts: it is not in $E_d(\mathcal{S})$, so it cannot be a TP or FP, and it does not match any element of $E(G_\pi)$, so the corresponding true edge falls into FN. Submitting unresolved orientations therefore reduces directed recall but never inflates directed false positives.

\subsection{Structural Hamming distance}
\label{app:scoring:shd}

SHD is computed over the union of true and submitted skeleton pairs,
\begin{equation}
\mathcal{P} = \pi(E(G_\pi)) \cup \mathrm{skel}(\mathcal{S}).
\end{equation}
For each $\{u, v\} \in \mathcal{P}$, exactly one of the following holds:
\begin{itemize}
\itemsep0pt
\item \textbf{Extra}: $\{u, v\} \notin \pi(E(G_\pi))$ (the submission asserts an adjacency that is not in the true DAG).
\item \textbf{Missing}: $\{u, v\} \in \pi(E(G_\pi))$ and $\{u, v\} \notin \mathrm{skel}(\mathcal{S})$.
\item \textbf{Unresolved}: $\{u, v\} \in E_u$ and $\{u, v\} \in \pi(E(G_\pi))$.
\item \textbf{Reversed}: $\{u, v\} \in \pi(E_d(\mathcal{S}))$ and $\{u, v\} \in \pi(E(G_\pi))$, but the submitted orientation is opposite to the true one.
\item \textbf{Correct}: $\{u, v\} \in \pi(E_d(\mathcal{S}))$, $\{u, v\} \in \pi(E(G_\pi))$, and the orientations match.
\end{itemize}
Each non-correct case contributes one error. The SHD is the total error count
\begin{equation}
\mathrm{SHD}(\mathcal{S}, G_\pi) = n_{\text{extra}} + n_{\text{missing}} + n_{\text{unresolved}} + n_{\text{reversed}},
\end{equation}
which is the four-error decomposition reported alongside SHD in the per-model tables.

\subsection{Worked example}
\label{app:scoring:example}

Take the 4-node DAG $G$ with edges $E(G) = \{1{\to}3,\; 2{\to}3,\; 1{\to}4\}$. Its CPDAG has the v-structure $\{1{\to}3,\; 2{\to}3\}$ as compelled edges and the edge $\{1, 4\}$ undirected: Meek rule~1 cannot orient $1 - 4$ because node $1$ has no parents to propagate from, and rules~2 and~3 require directed edges incident to node~$4$, of which there are none. Consider a submission $\mathcal{S}$ with directed edges $E_d(\mathcal{S}) = \{(3, 1),\; (1, 4),\; (2, 4)\}$ and undirected edges $E_u(\mathcal{S}) = \{\{2, 3\}\}$. The metrics evaluate as follows.

\begin{center}
\begin{tabular}{lll}
\toprule
Metric              & Counts                              & Value \\
\midrule
Skeleton F1         & TP $= 3$, FP $= 1$, FN $= 0$        & $\mathrm{F1} = 6/7 \approx 0.857$ \\
Compelled-edge F1   & TP $= 0$, FP $= 1$, FN $= 2$        & $\mathrm{F1} = 0$ \\
Directed F1         & TP $= 1$, FP $= 2$, FN $= 2$        & $\mathrm{F1} = 1/3$ \\
SHD                 & $1$ reversed, $1$ unresolved, $1$ extra & $\mathrm{SHD} = 3$ \\
\bottomrule
\end{tabular}
\end{center}

The compelled-edge edge $(1, 3)$ is submitted reversed as $(3, 1)$ and contributes one FP and one FN. The compelled-edge edge $(2, 3)$ is submitted as unresolved and contributes one FN, no FP. The directed-F1 numerator credits only $(1, 4)$, which appears in both $E_d(\mathcal{S})$ and $E(G)$. The SHD pair $(1, 3)$ is reversed, $(2, 3)$ is unresolved, $(2, 4)$ is extra, and $(1, 4)$ is correct.
